\section{The Spirit of Joseph de Maistre}

\textbf{Joseph de Maistre} was an Hermetist (in the school of \textbf{Martines de Pasqually} and \textbf{Louis-Claude de Saint-Martin}), Traditionalist, Reactionary, European, and a source of the Russian Idea. In his Self-Defense, Julius Evola includes Maistre among those well-bred men holding sane and normal views common prior to the French Revolution. In this brief review of the Maistre's St Petersburg Dialogues Evola again reveals his high regard for de Maistre, despite their fundamental disagreement on religion and providence. 

In \textit{Materialism and the Task of Anthroposophy}, Rudolf Steiner says this about Joseph de Maistre:

\begin{quotex}
With de Maistre, a personality of the greatest imaginable genius, of compelling spirituality but Roman Catholic through and through, appears on the scene.

\end{quotex}
High praise indeed, especially from a man who was in fundamental disagreement with de Maistre.

\paragraph{Digression} To complete the circle, we have mentioned that Evola included Steiner's picture not once, but twice, as demonstrating ``spiritual penetration", despite being in fundamental disagreement with Steiner's doctrines of Karma and Reincarnation. So here we have the example of three esoterists demonstrating respect despite their disagreements on specific issues. This is how men of high intellectual attainment behave, since they understand what ``thinking" is. Lesser men treat thought as a weapon, and means of combat. They confuse the use of words and ideas, with the words and ideas themselves. Let this serve as a model.

\paragraph{Returning to the topic} Steiner has some interesting things to say about de Masitre who represents the reactionary trend to Steiner's progressivist view of mankind. This is how he characterizes de Maistre's basic worldview:

\begin{quotex}
All of humanity falls into two categories, one representing the kingdom of God, the other representing the kingdom of this world. 

\end{quotex}
As we recently pointed out, this is also the view of \textbf{Vladimir Soloviev} (not to mention \textbf{St. Augustine}), and the reason for this seeming coincidence will become clear. But the point is that there is an essential difference in the quality of being between the two groups of humanity: one spiritually aware, the other not. This is more than a superficial adherence to one system of belief over another, but represents a real change in the level of being.

Steiner writes:

\begin{quotex}
there are those in Europe who cling to this view that since the beginning of the fifteenth century divine world rule has assumed a quite different position in regard to earth humanity. 

\end{quotex}
Of course, this is not just representative of the Catholicism of de Maistre, but it also reflects the point of view of both Julius Evola and Rene Guenon, who saw the Reformation as the beginning of the end of the last vestiges of Traditional civilization in the West.

Steiner astutely points this out:

\begin{quotex}
De Maistre had the grandiose idea to tie in with Russianism, particularly with the element that had found its way since ancient times from Asia into the Orthodox Catholic, Russian religion. From there, he wished to create a connection to Romanism. He tried to bring about the great fusion between the element living in the Oriental manner of thinking in Russian culture, and the element coming from Rome. already imbued with this view. 

\end{quotex}
Here we see this revealed in the writings Soloviev who in his own way tries to tie the ``Oriental manner of thinking" into contemporary Western ideas.

Steiner continues:

\begin{quotex}
de Maistre refers back to what Christianity was in regard to its metaphysical view prior to the scholastic age, what it was in the first centuries and what was acceptable to Rome. De Maistre aimed for Roman, for Catholic, Christianity as a real power.

\end{quotex}
Again, as Evola also points out, the Catholic counter-reformation was not enough, it only went back half-way. It needed to return to the primordial sources of its spiritual power. Instead of being the opposite of a Reformation, it merely contented itself with being a reformation in the opposite direction. The spiritually weak state of Christianity today is the result.

Steiner's short essay puts the issue in perfect clarity.

\begin{itemize}
\item Is the human race on the way to some progressive, evolutionary, utopian future? 
\item Or is there a fundamental and irreconcilable divide between those on the side of Tradition and those on the side of Revolution? 
\end{itemize}
If the latter\footnote{The idea the evolution and progress of the human race is one of the features that distinguish ``New Age" teachings from Tradition.}, then the task is to begin the reconstruction of Tradition from the insights given by de Maistre, Soloviev, Guenon, and Evola.


\flrightit{Posted on 2010-05-23 by Cologero }
